% =====================================================================
%  Chương 1 — Hạng mục 1: cơ sở lý thuyết mạng nơ-ron hồi quy
%  Mã nguồn: a06/theory.py · Notebook: 00_ly_thuyet_rnn.ipynb
% =====================================================================
\chapter{Cơ sở lý thuyết mạng nơ-ron hồi quy}
\label{ch:lythuyet}

Chương này dựng lại các khái niệm của mạng hồi quy từ công thức. Mỗi khái niệm là một hàm NumPy trong
\texttt{a06/theory.py}; notebook 00 chạy hàm đó và so với lớp PyTorch tương ứng bằng \texttt{np.allclose}. Cả
\SL{theory/n_checks} phép đối chiếu đều cho kết quả khớp, nên các công thức dưới đây chính là những gì thư viện
đang tính.

\section{Dữ liệu chuỗi và tensor ba chiều}

Một mẫu dữ liệu chuỗi là một dãy vector $x_1, x_2, \dots, x_T$ với $x_t \in \mathbb{R}^D$. Gom $B$ mẫu lại được
tensor $X \in \mathbb{R}^{B \times T \times D}$ (Hình~\ref{fig:tensor}). Cả hai tập của bài đều có dạng này. AMZN
cho $X_{\text{train}}$ có shape $(\SL{theory/shapes/amzn/0}, \SL{theory/shapes/amzn/1}, \SL{theory/shapes/amzn/2})$:
mỗi mẫu là một cửa sổ 30 phiên giao dịch, mỗi phiên 8 đặc trưng. KKBox cho
$(\SL{theory/shapes/kkbox/0}, \SL{theory/shapes/kkbox/1}, \SL{theory/shapes/kkbox/2})$: mỗi mẫu là một người dùng,
31 ngày của tháng 3/2017, mỗi ngày 8 đặc trưng nghe nhạc.

\begin{figure}[H]
  \centering
  \hinh{f1_tensor}
  \caption{Tensor dữ liệu chuỗi; lát tô cam là một bước thời gian $x_t$}
  \label{fig:tensor}
\end{figure}

Điều làm dữ liệu chuỗi khác dữ liệu bảng là thứ tự. Giá hôm nay phụ thuộc giá hôm qua; một người sắp bỏ dịch vụ
thường nghe ít dần trong những ngày cuối. Giả định các mẫu độc lập và cùng phân phối (i.i.d.) vẫn đúng giữa các
mẫu, nhưng sai hoàn toàn giữa các bước trong cùng một mẫu. Mô hình cho chuỗi phải khai thác được sự phụ thuộc đó.

\section{Vì sao MLP và Conv1D chưa đủ}

Có hai cách dùng lại công cụ của các bài trước. Cách thứ nhất là làm phẳng chuỗi thành vector $T \cdot D$ chiều rồi
đưa vào MLP. Cách thứ hai là trượt một kernel Conv1D theo trục thời gian. Bảng~\ref{tab:mlp-rnn} đếm tham số của tầng
đầu tiên với $T = 30$, $D = 8$ và $H = 64$ đơn vị ẩn.

\begin{table}[H]
  \centering\small
  \caption{Tầng đầu của ba loại mô hình cho chuỗi $T = 30$, $D = 8$, $H = 64$}
  \label{tab:mlp-rnn}
  \begin{tabular}{lL{3.3cm}rL{3.2cm}L{3.2cm}}
    \toprule
    Mô hình & Số tham số & Giá trị & Độ dài chuỗi thay đổi & Nhớ quá khứ xa \\
    \midrule
    MLP (làm phẳng) & $(T \cdot D) H + H$ & \SL{theory/params/mlp} & không, $T$ cố định & mọi bước, nhưng không chia sẻ trọng số \\
    Conv1D, $K = 3$ & $K D H + H$ & \SL{theory/params/conv1d} & được & chỉ $K$ bước mỗi tầng \\
    Simple RNN & $H(D + H + 1)$ & \SL{theory/params/rnn} & được & qua trạng thái $h_t$ \\
    \bottomrule
  \end{tabular}
\end{table}

MLP học một trọng số riêng cho từng cặp (bước, đặc trưng). Nhân đôi độ dài chuỗi thì số tham số tầng đầu tăng từ
\SL{theory/params/mlp} lên \SL{theory/params/mlp_t60}, và điều học được ở bước 5 không dùng lại được ở bước 25.
Conv1D chia sẻ trọng số theo thời gian nhưng mỗi tầng chỉ nhìn $K$ bước, muốn thấy 30 bước phải chồng nhiều tầng.
RNN giải quyết cả hai: một bộ trọng số dùng cho mọi bước, và một trạng thái $h_t$ mang thông tin từ đầu chuỗi đến
hiện tại. Ba hàm đếm tham số nằm trong \texttt{theory.py}; notebook so từng công thức với \texttt{numel()} của
\texttt{nn.Linear}, \texttt{nn.Conv1d} và \texttt{nn.RNN}.

\maNguon{theory__count_params_rnn}{Số tham số của Simple RNN không phụ thuộc độ dài chuỗi}

\section{Simple RNN}

Elman~\cite{elman1990} đưa ra dạng đơn giản nhất: ở mỗi bước, trạng thái mới là hàm của đầu vào hiện tại và trạng
thái cũ,
\begin{equation}
  h_t = \tanh\left(W_{xh}\, x_t + W_{hh}\, h_{t-1} + b\right), \qquad h_0 = 0 .
  \label{eq:rnn}
\end{equation}
Mở vòng lặp theo thời gian thì RNN trở thành một mạng sâu $T$ tầng, nhưng mọi tầng dùng chung $W_{xh}$, $W_{hh}$ và
$b$ (Hình~\ref{fig:unroll}). Trạng thái cuối $h_T$ là bản tóm tắt cả chuỗi, và cả bốn mô hình của bài đều đưa
$h_T$ vào một đầu dự báo nhỏ.

\begin{figure}[H]
  \centering
  \hinh{f1_unroll}
  \caption{RNN ở dạng cuộn và dạng mở theo thời gian}
  \label{fig:unroll}
\end{figure}

\maNguon{theory__rnn_forward}{Lan truyền xuôi của Simple RNN theo công thức~\eqref{eq:rnn}}

Hàm trên khớp \texttt{nn.RNN} của PyTorch (kiểm tra \texttt{rnn\_forward}: \SL{theory/checks/rnn_forward}). PyTorch lưu
hai vector bias $b_{ih}$ và $b_{hh}$ nên tầng \texttt{nn.RNN(8, 64)} có \SL{theory/cell_params/rnn} tham số, nhiều hơn
công thức ở Bảng~\ref{tab:mlp-rnn} đúng $H = 64$. Khi so khớp, notebook cộng hai bias lại thành $b$.

\section{Lan truyền ngược theo thời gian và hai vấn đề của gradient}

Huấn luyện RNN dùng lan truyền ngược theo thời gian (BPTT)~\cite{werbos1990}: mở mạng ra $T$ bước rồi lan truyền
ngược như một mạng thường, cộng gradient của $W_{hh}$ từ mọi bước. Gọi $L$ là hàm mất mát tính từ $h_T$. Theo quy tắc
chuỗi,
\begin{equation}
  \frac{\partial L}{\partial h_t}
  = \frac{\partial L}{\partial h_T} \prod_{k=t+1}^{T} \frac{\partial h_k}{\partial h_{k-1}}
  = \frac{\partial L}{\partial h_T} \prod_{k=t+1}^{T} W_{hh}^\top\, \mathrm{diag}\!\left(1 - h_k^2\right).
  \label{eq:bptt}
\end{equation}
Gradient đi từ $h_T$ về $h_t$ bị nhân với $T - t$ ma trận Jacobi. Nếu chuẩn của các ma trận này nhỏ hơn 1, tích co lại
theo hàm mũ và gradient \emph{triệt tiêu}; nếu lớn hơn 1, tích \emph{bùng nổ}~\cite{bengio1994,pascanu2013}. Bỏ qua
đạo hàm của $\tanh$ (luôn không vượt 1) để thấy riêng vai trò của $W_{hh}$: với vector $v$ bất kỳ,
$\lVert (W_{hh}^\top)^t v \rVert$ tăng hay giảm theo bán kính phổ $\rho$ của $W_{hh}$.

\maNguon{theory__jacobian_product_norm}{Chuẩn của tích Jacobi với bán kính phổ $\rho$ cho trước}

Hình~\ref{fig:rho} chạy hàm này với ba giá trị $\rho$. Sau \SL{theory/rho_end/T} bước, $\rho = 0{,}9$ còn lại
\SL{theory/rho_end/r09}, còn $\rho = 1{,}1$ đã lên \SL{theory/rho_end/r11}. Chênh lệch chỉ $0{,}1$ ở bán kính phổ
thành chênh lệch hơn ba bậc độ lớn ở gradient.

\begin{figure}[H]
  \centering
  \hinh{f1_rho}
  \caption{Tích Jacobi co hoặc giãn theo hàm mũ tuỳ bán kính phổ của $W_{hh}$}
  \label{fig:rho}
\end{figure}

Thực tế có thêm $\tanh$ và trọng số khởi tạo ngẫu nhiên, nên notebook đo trực tiếp: dựng chuỗi ngẫu nhiên dài
\SL{theory/grad/T} bước, lấy $L = \sum h_T$, rồi dùng autograd đọc $\lVert\partial L/\partial h_t\rVert$ ở từng bước.
Để giữ được gradient của từng $h_t$, hàm mở vòng lặp bằng \texttt{RNNCell}/\texttt{LSTMCell} thay cho lớp RNN gộp.

\maNguon{theory__grad_norms_through_time}{Đo chuẩn gradient theo từng bước bằng autograd}

\begin{figure}[H]
  \centering
  \hinh{f1_gradnorm}
  \caption{Gradient của mất mát theo $h_t$, chuẩn hoá theo bước cuối (trục tung thang log)}
  \label{fig:gradnorm}
\end{figure}

Hình~\ref{fig:gradnorm} cho một kết quả đáng chú ý. Với Simple RNN, lùi 20 bước thì gradient chỉ còn
\SL{theory/grad/rnn_lag20} lần giá trị ở bước cuối; lùi 59 bước còn \SL{theory/grad/rnn_lag59}. LSTM với khởi tạo
mặc định của PyTorch cũng suy giảm gần như vậy: \SL{theory/grad/lstm_lag20} ở trễ 20. Lý do là bias cổng quên bằng 0
nên $f_t \approx 0{,}5$, và cell state bị nhân với khoảng $0{,}5$ ở mỗi bước. Khi đặt bias cổng quên bằng
\SL{theory/grad/forget_bias} như Jozefowicz và cộng sự đề xuất~\cite{jozefowicz2015}, $f_t$ tiến về gần 1 và gradient ở trễ
59 còn \SL{theory/grad/lstm3_lag59}, lớn hơn RNN hàng chục bậc. Vậy LSTM không tự động nhớ xa. Nó có một con đường để
gradient đi qua, và mô hình phải học mở con đường đó.

\section{Gradient clipping}

Triệt tiêu làm mô hình không học được phụ thuộc xa; bùng nổ thì nguy hiểm hơn, vì một bước cập nhật quá lớn có thể phá
hỏng trọng số đã học. Cách chữa thông dụng là cắt theo chuẩn~\cite{pascanu2013}: gộp mọi gradient thành một vector
$g$, nếu $\lVert g\rVert > c$ thì thay bằng $g \cdot c / \lVert g\rVert$. Độ dài bị chặn, còn hướng giữ nguyên
(Hình~\ref{fig:clip}). Cả 16 mô hình của bài dùng $c = 1$.

\maNguon{theory__clip_grad_norm}{Cắt gradient theo chuẩn, giữ hướng}

\begin{figure}[H]
  \centering
  \hinh{f1_clip}
  \caption{Gradient clipping chỉ đổi độ dài của $g$, không đổi hướng}
  \label{fig:clip}
\end{figure}

Trên hai ma trận gradient ngẫu nhiên có tổng chuẩn \SL{theory/clip/norm_before}, hàm trả về chuẩn
\SL{theory/clip/norm_after} và khớp \texttt{nn.utils.clip\_grad\_norm\_} (\SL{theory/checks/clip_grad_norm}).

\section{LSTM}

LSTM~\cite{hochreiter1997} thêm một trạng thái thứ hai là cell state $c_t$ và ba cổng điều khiển nó. Với
$z_t = [h_{t-1}; x_t]$,
\begin{align}
  f_t &= \sigma(W_f z_t + b_f), & i_t &= \sigma(W_i z_t + b_i), & o_t &= \sigma(W_o z_t + b_o), \nonumber\\
  \tilde c_t &= \tanh(W_c z_t + b_c), & c_t &= f_t \odot c_{t-1} + i_t \odot \tilde c_t, & h_t &= o_t \odot \tanh(c_t).
  \label{eq:lstm}
\end{align}
Cổng quên $f_t$ quyết định giữ bao nhiêu của $c_{t-1}$, cổng vào $i_t$ quyết định ghi bao nhiêu thông tin mới
$\tilde c_t$, cổng ra $o_t$ quyết định để lộ bao nhiêu ra $h_t$. Điểm mấu chốt là phép cộng ở dòng $c_t$
(Hình~\ref{fig:lstm}): $\partial c_t / \partial c_{t-1} = f_t$, không có ma trận trọng số và không có $\tanh$ trên
đường đi. Khi $f_t$ gần 1, gradient đi qua nhiều bước gần như nguyên vẹn.

\begin{figure}[H]
  \centering
  \hinh{f1_lstm}
  \caption{Tế bào LSTM: đường cộng của cell state ở trên, bốn phép biến đổi có cổng ở dưới}
  \label{fig:lstm}
\end{figure}

\maNguon{theory__lstm_cell}{Một bước LSTM theo~\eqref{eq:lstm}, thứ tự khối cổng như PyTorch}

\section{GRU}

GRU~\cite{cho2014} gộp cell state vào $h_t$ và chỉ dùng hai cổng: cổng reset $r_t$ quyết định dùng bao nhiêu trạng thái
cũ khi tạo ứng viên $n_t$, cổng cập nhật $z_t$ hoà trộn trạng thái cũ với ứng viên,
\begin{align}
  r_t &= \sigma(W_r z_t + b_r), \qquad z_t = \sigma(W_z z_t + b_z), \nonumber\\
  n_t &= \tanh\!\left(W_{in} x_t + r_t \odot (W_{hn} h_{t-1})\right), \qquad
  h_t = (1 - z_t) \odot n_t + z_t \odot h_{t-1}.
  \label{eq:gru}
\end{align}
Dòng cuối là một tổ hợp lồi: khi $z_t \to 1$ thì $h_t \approx h_{t-1}$, gradient đi thẳng qua giống đường cộng của
LSTM (Hình~\ref{fig:gru}).

\begin{figure}[H]
  \centering
  \hinh{f1_gru}
  \caption{Tế bào GRU với hai cổng}
  \label{fig:gru}
\end{figure}

\maNguon{theory__gru_cell}{Một bước GRU theo~\eqref{eq:gru}}

Bảng~\ref{tab:lstm-gru} so hai tế bào với cùng $D = 8$, $H = 64$. Số tham số đọc từ PyTorch.

\begin{table}[H]
  \centering\small
  \caption{So sánh Simple RNN, LSTM, GRU và BiLSTM ($D = 8$, $H = 64$, tham số của riêng tầng hồi quy trong PyTorch)}
  \label{tab:lstm-gru}
  \begin{tabular}{lcccr}
    \toprule
    Tế bào & Số khối trọng số & Trạng thái & Cổng & Tham số \\
    \midrule
    Simple RNN & 1 & $h_t$ & không & \SL{theory/cell_params/rnn} \\
    LSTM & 4 & $h_t$, $c_t$ & $f, i, o$ & \SL{theory/cell_params/lstm} \\
    GRU & 3 & $h_t$ & $r, z$ & \SL{theory/cell_params/gru} \\
    BiLSTM & $2 \times 4$ & $\overrightarrow{h}_t, \overleftarrow{h}_t$ & $f, i, o$ mỗi chiều & \SL{theory/cell_params/bilstm} \\
    \bottomrule
  \end{tabular}
\end{table}

\section{LSTM hai chiều}

BiLSTM~\cite{schuster1997} chạy hai LSTM độc lập, một đọc từ $t = 1$ đến $T$, một đọc từ $T$ về 1, rồi nối hai trạng
thái cuối: $[\overrightarrow{h}_T; \overleftarrow{h}_1] \in \mathbb{R}^{2H}$ (Hình~\ref{fig:bilstm}). Notebook kiểm tra
cách nối này trên \texttt{nn.LSTM(bidirectional=True)}: đầu ra có shape
$(\SL{theory/bilstm/out_shape/0}, \SL{theory/bilstm/out_shape/1}, \SL{theory/bilstm/out_shape/2})$, nửa đầu ở bước cuối
trùng $\overrightarrow{h}_T$ và nửa sau ở bước đầu trùng $\overleftarrow{h}_1$ (\SL{theory/checks/bilstm_concat}).

\begin{figure}[H]
  \centering
  \hinh{f1_bilstm}
  \caption{BiLSTM: vector tóm tắt là trạng thái cuối của hai chiều đọc}
  \label{fig:bilstm}
\end{figure}

Với bài dự báo giá, chiều ngược không làm lộ tương lai. Nó chỉ đọc lại 30 phiên của cửa sổ theo thứ tự ngược, mọi phiên
đều trước ngày cần dự báo. Chiều ngược thêm một góc nhìn khác lên cùng quá khứ, không thêm thông tin mới.

\section{Điều chuẩn cho mạng hồi quy}

Dropout thường đổi mặt nạ ở mỗi bước. Áp lên trạng thái $h_t$, nó xoá một phần trí nhớ khác nhau ở mỗi bước, và
thông tin khó đi xa. Gal và Ghahramani~\cite{gal2016} đề xuất dùng \emph{một} mặt nạ cho cả chuỗi (recurrent dropout):
đơn vị bị tắt thì tắt suốt chuỗi, đơn vị còn lại mang thông tin liền mạch (Hình~\ref{fig:dropout}, $p =
\SL{theory/mask_p}\,\%$).

\begin{figure}[H]
  \centering
  \hinh{f1_dropout}
  \caption{Mặt nạ dropout thường và recurrent dropout; ô đỏ là đơn vị bị tắt}
  \label{fig:dropout}
\end{figure}

\maNguon{theory__recurrent_dropout_masks}{Hai kiểu mặt nạ dropout cho chuỗi}

BatchNorm của các bài trước lấy thống kê theo lô và theo từng bước, không hợp với chuỗi có độ dài thay đổi.
LayerNorm~\cite{ba2016} chuẩn hoá theo chiều đặc trưng của \emph{từng mẫu}, nên dùng được ở mọi bước và mọi kích thước lô.
Hàm dưới khớp \texttt{nn.LayerNorm} (\SL{theory/checks/layer_norm}).

\maNguon{theory__layer_norm}{Layer Normalization theo từng mẫu}

Các mô hình ở Chương~\ref{ch:pytorch} chỉ có một tầng hồi quy 64 đơn vị và dừng sớm theo validation, nên bài không
dùng dropout hay LayerNorm. Hai kỹ thuật được trình bày ở đây vì chúng là công cụ chuẩn khi mô hình sâu hơn.
